% !TEX program = xelatex
% =============================================================================
%  墨经 Mojing · 数学论文模板示例
%  Bernstein 多项式逼近的收敛阶
%
%  编译：latexmk main.tex   （自动调用 biber 处理参考文献）
%  图表数据由 data/make_data.py 生成，可直接重新运行以核对数值。
%  选项：preprint（预印本样式，不排刊头）、thmstyle=song（定理正文用宋体）、
%        enabstract=front（英文摘要排在首页，默认排在文末）、oneside
% =============================================================================
\documentclass{mojing-paper}
\addbibresource{refs.bib}

\journalinfo{
  name      = 墨经数学通讯,
  enname    = Mojing Mathematical Letters,
  volume    = 1, issue = 1, year = 2026, month = 10,
  doi       = 10.0000/mojing.2026.0001,
  arxiv     = 2610.00000,   % 预印本样式（preprint 选项）时显示
  firstpage = 1,
}

\title[Bernstein 多项式逼近的收敛阶]
  {Bernstein 多项式逼近的收敛阶：\\一个概率论证明及其精确性}
\author[张三，李四]{张\quad 三\affil{1}，\quad 李\quad 四\affil{1,2}}
\affiliation{1.\ 示例大学数学科学学院，北京 100000；\quad
             2.\ 示例大学统计与数据科学研究中心，北京 100000}
\entitle{The Rate of Convergence of Bernstein Polynomials:\\
         A Probabilistic Proof and Its Sharpness}
\enauthor{ZHANG San\textsuperscript{1}, \ LI Si\textsuperscript{1,2}}
\enaffiliation{1. School of Mathematical Sciences, Example University, Beijing 100000, China;\\
               2. Center for Statistics and Data Science, Example University, Beijing 100000, China}

\begin{abstract}
本文用概率论的语言给出 Popoviciu 关于 Bernstein 多项式逼近阶的估计
$\norm{B_nf-f}\le\frac32\,\omega(f,n^{-1/2})$ 的一个简短证明。通过对函数
$f_0(x)=\abs{x-\frac12}$ 的精确计算，我们证明
$B_{2m}f_0(\frac12)=\binom{2m}{m}\big/\bigl(2\cdot4^m\bigr)\sim(2\pi n)^{-1/2}$，从而阶 $n^{-1/2}$
对 Lipschitz 函数类不能改进。对二阶连续可微函数，我们给出 Voronovskaya 渐近公式的一个
定量形式：$\abs{n(B_nf-f)-\frac12x(1-x)f''}\le\frac{7}{32}\,\omega(f'',n^{-1/2})$。
数值实验印证了上述结论。
\end{abstract}
\keywords{Bernstein 多项式；连续模；收敛阶；Voronovskaya 公式；二项分布}
\msc{41A10；41A25；41A36}

\begin{enabstract}
We give a short probabilistic proof of Popoviciu's estimate
$\norm{B_nf-f}\le\frac32\,\omega(f,n^{-1/2})$ for Bernstein polynomials. An exact computation for
$f_0(x)=\abs{x-\frac12}$ shows that $B_{2m}f_0(\frac12)=\binom{2m}{m}/(2\cdot4^m)\sim(2\pi n)^{-1/2}$,
so the order $n^{-1/2}$ cannot be improved on the class of Lipschitz functions. For
$f\in C^2[0,1]$ we prove a quantitative Voronovskaya formula with the error bound
$\frac{7}{32}\,\omega(f'',n^{-1/2})$. Numerical experiments illustrate the results.
\end{enabstract}
\enkeywords{Bernstein polynomials; modulus of continuity; rate of convergence;
Voronovskaya formula; binomial distribution}

\received{2026-06-18}
\revised{2026-09-02}
\funding{示例基金资助项目（编号：0000000）。本文为排版模板示例，作者与单位均为虚构。}
\authorinfo{张三（1990—），男，博士，主要研究方向为函数逼近论。
通信作者：李四，E-mail：lisi@example.edu.cn}

\begin{document}
\maketitle

% =============================================================================
\section{引言}
% =============================================================================

Weierstrass 于 1885 年证明，闭区间上的连续函数可以用多项式一致逼近\cite{weierstrass1885}。
1912 年，Bernstein 给出了一个构造性的、带有鲜明概率色彩的证明\cite{bernstein1912}：对
$f\in C[0,1]$，令
\begin{equation}\label{eq:bernstein}
  B_nf(x)=\sum_{k=0}^{n}f\Bigl(\frac{k}{n}\Bigr)\binom{n}{k}x^k(1-x)^{n-k},\qquad x\in[0,1],
\end{equation}
则 $B_nf$ 在 $[0,1]$ 上一致收敛于 $f$。若 $S_n$ 服从参数为 $n$、$x$ 的二项分布，记作
$S_n\sim\mathrm{Bin}(n,x)$，则 \cref{eq:bernstein} 可以写成
\begin{equation}\label{eq:expectation}
  B_nf(x)=\mathbb{E}\,f\Bigl(\frac{S_n}{n}\Bigr).
\end{equation}
由大数定律，$S_n/n$ 依概率收敛于 $x$，而 $f$ 一致连续，于是 $B_nf(x)\to f(x)$。这正是
Bernstein 原始证明的思路。

本文关心收敛的速度。记 $\norm{\cdot}$ 为 $[0,1]$ 上的上确界范数，
\[
  \omega(f,\delta)=\sup\set{\abs{f(s)-f(t)}\given s,t\in[0,1],\ \abs{s-t}\le\delta}
\]
为 $f$ 的连续模。Popoviciu\cite{popoviciu1935}证明了
$\norm{B_nf-f}\le\frac32\,\omega(f,n^{-1/2})$；Sikkema\cite{sikkema1961}进一步求出了其中的最优常数。
对于光滑函数，Voronovskaya\cite{voronovskaya1932}得到了渐近公式
$\lim_{n\to\infty}n\bigl(B_nf(x)-f(x)\bigr)=\frac12x(1-x)f''(x)$。这些结果的系统论述可参见专著
\cite{lorentz1953,devore1993}。

本文的目的是在统一的概率框架下处理上述结果。我们的主要贡献有三点：
\begin{enumerate}
  \item 利用 \cref{eq:expectation} 与二阶矩，给出 Popoviciu 估计的一个四行证明（\cref{thm:popoviciu}）；
  \item 通过精确计算说明阶 $n^{-1/2}$ 对 Lipschitz 函数类是最优的，并给出误差的渐近常数
        $(2\pi)^{-1/2}$（\cref{thm:sharp}）；
  \item 利用四阶矩，得到 Voronovskaya 公式的一个带显式常数的定量形式（\cref{thm:voronovskaya}）。
\end{enumerate}
第 2 节准备二项分布的矩估计，第 3--5 节依次证明上述结果，第 6 节为结语。

% =============================================================================
\section{二项分布的矩}
% =============================================================================

以下总设 $S_n\sim\mathrm{Bin}(n,x)$，$x\in[0,1]$，并记 $T_n=S_n/n-x$。

\begin{lemma}\label{lem:moments}
对一切 $n\ge1$ 与 $x\in[0,1]$，有
\begin{equation}\label{eq:moments}
  \mathbb{E}\,T_n=0,\qquad
  \mathbb{E}\,T_n^2=\frac{x(1-x)}{n}\le\frac{1}{4n},\qquad
  \mathbb{E}\,T_n^4=\frac{x(1-x)\bigl[1+3(n-2)x(1-x)\bigr]}{n^3}<\frac{3}{16n^2}.
\end{equation}
\end{lemma}

\begin{proof}
把 $S_n$ 写成 $n$ 个独立的 Bernoulli 变量之和 $S_n=\sum_{i=1}^n\xi_i$，$\mathbb{P}(\xi_i=1)=x$。
令 $\eta_i=\xi_i-x$，则 $\mathbb{E}\,\eta_i=0$，$\mathbb{E}\,\eta_i^2=x(1-x)$，
$\mathbb{E}\,\eta_i^4=x(1-x)\bigl(1-3x(1-x)\bigr)$。展开 $(\sum\eta_i)^4$ 并利用独立性，只有
$\eta_i^4$ 与 $\eta_i^2\eta_j^2$（$i\ne j$）两类项的期望非零，故
\[
  \mathbb{E}\Bigl(\sum_{i=1}^n\eta_i\Bigr)^4=n\,\mathbb{E}\,\eta_1^4+3n(n-1)\bigl(\mathbb{E}\,\eta_1^2\bigr)^2
  =nx(1-x)\bigl[1+3(n-2)x(1-x)\bigr].
\]
除以 $n^4$ 即得 \cref{eq:moments} 中的等式（二阶矩同理）。最后，由 $x(1-x)\le\frac14$，
$\mathbb{E}\,T_n^4\le\frac{1}{4n^3}\bigl(1+\frac{3(n-2)}{4}\bigr)=\frac{3n-2}{16n^3}<\frac{3}{16n^2}$。
\end{proof}

我们还需要连续模的一个熟知性质，其证明可见文献 \parencite[第 2 章]{devore1993}。

\begin{lemma}\label{lem:modulus}
对 $f\in C[0,1]$ 及 $\delta>0$，$h\ge0$，有
$\omega(f,h)\le\bigl(1+\frac{h}{\delta}\bigr)\omega(f,\delta)$，并且
$\omega(f,h)\le\bigl(1+\frac{h^2}{\delta^2}\bigr)\omega(f,\delta)$。
\end{lemma}

\begin{proof}
第一式来自次可加性 $\omega(f,h_1+h_2)\le\omega(f,h_1)+\omega(f,h_2)$：
$\omega(f,h)\le\lceil h/\delta\rceil\,\omega(f,\delta)\le(1+h/\delta)\,\omega(f,\delta)$。
对第二式，$h\le\delta$ 时由单调性显然；$h>\delta$ 时 $h/\delta<h^2/\delta^2$，由第一式即得。
\end{proof}

% =============================================================================
\section{Popoviciu 估计}
% =============================================================================

\begin{theorem}[Popoviciu]\label{thm:popoviciu}
对一切 $f\in C[0,1]$ 与 $n\ge1$，有
\[
  \norm{B_nf-f}\le\frac32\,\omega\Bigl(f,\frac{1}{\sqrt n}\Bigr).
\]
\end{theorem}

\begin{proof}
固定 $x\in[0,1]$ 与 $\delta>0$。由 \cref{eq:expectation}、\cref{lem:modulus} 与 Cauchy--Schwarz 不等式，
\[
  \abs{B_nf(x)-f(x)}\le\mathbb{E}\,\abs{f(x+T_n)-f(x)}
  \le\omega(f,\delta)\Bigl(1+\frac{\mathbb{E}\abs{T_n}}{\delta}\Bigr)
  \le\omega(f,\delta)\Bigl(1+\frac{(\mathbb{E}\,T_n^2)^{1/2}}{\delta}\Bigr).
\]
由 \cref{lem:moments}，$(\mathbb{E}\,T_n^2)^{1/2}\le\frac{1}{2\sqrt n}$。取 $\delta=n^{-1/2}$ 即得。
\end{proof}

\begin{corollary}\label{cor:lipschitz}
若 $f$ 满足 Lipschitz 条件 $\abs{f(s)-f(t)}\le L\abs{s-t}$，则 $\norm{B_nf-f}\le\dfrac{L}{2\sqrt n}$。
\end{corollary}

\begin{proof}
$\abs{B_nf(x)-f(x)}\le L\,\mathbb{E}\abs{T_n}\le L\,(\mathbb{E}\,T_n^2)^{1/2}\le\frac{L}{2\sqrt n}$。
\end{proof}

\begin{remark}
\cref{thm:popoviciu} 中的常数 $\frac32$ 不是最优的。Sikkema\cite{sikkema1961}证明了最优常数为
$\frac{4306+837\sqrt6}{5832}\approx1.0899$。
\end{remark}

% =============================================================================
\section{阶的精确性}
% =============================================================================

自然要问：阶 $n^{-1/2}$ 能否改进？下面的定理表明，对 Lipschitz 函数类，答案是否定的，
而且 \cref{cor:lipschitz} 中的常数 $\frac12$ 至多只能改进到 $(2\pi)^{-1/2}\approx0.3989$。

\begin{theorem}\label{thm:sharp}
设 $f_0(x)=\abs{x-\frac12}$，则对一切正整数 $m$，
\begin{equation}\label{eq:exact}
  B_{2m}f_0\Bigl(\frac12\Bigr)=\frac{1}{2\cdot4^m}\binom{2m}{m}.
\end{equation}
特别地，当 $n=2m\to\infty$ 时，$B_nf_0(\frac12)-f_0(\frac12)\sim\dfrac{1}{\sqrt{2\pi n}}$。
\end{theorem}

\begin{proof}
取 $x=\frac12$，$S\sim\mathrm{Bin}(2m,\frac12)$。由 \cref{eq:expectation}，
\[
  B_{2m}f_0\Bigl(\frac12\Bigr)=\mathbb{E}\,\Bigl\lvert\frac{S}{2m}-\frac12\Bigr\rvert
  =\frac{1}{2m\cdot4^m}\sum_{k=0}^{2m}\abs{k-m}\binom{2m}{k}
  =\frac{1}{m\cdot4^m}\sum_{k=m+1}^{2m}(k-m)\binom{2m}{k},
\]
最后一步用了对称性 $\binom{2m}{k}=\binom{2m}{2m-k}$。利用
$k\binom{2m}{k}=2m\binom{2m-1}{k-1}$ 与 $\binom{2m}{k}=\binom{2m-1}{k}+\binom{2m-1}{k-1}$，得
\[
  (k-m)\binom{2m}{k}=m\left[\binom{2m-1}{k-1}-\binom{2m-1}{k}\right].
\]
对 $k=m+1,\dotsc,2m$ 求和，右端相消为 $m\binom{2m-1}{m}=\frac m2\binom{2m}{m}$，代入即得
\cref{eq:exact}。最后，由 Stirling 公式（见文献 \parencite{zhang1990,feller1968}），
$\binom{2m}{m}4^{-m}\sim(\pi m)^{-1/2}$，于是
$B_{2m}f_0(\frac12)\sim\frac{1}{2\sqrt{\pi m}}=\frac{1}{\sqrt{2\pi n}}$。
\end{proof}

\Cref{tab:sharp} 列出了 $e_n=B_nf_0(\frac12)$ 的精确值。可以看到 $\sqrt n\,e_n$ 自下而上
单调地趋于 $(2\pi)^{-1/2}=0.398942\dotso$，而 $e_n$ 与 Popoviciu 上界之比稳定在
$\frac{2}{3\sqrt{2\pi}}\approx0.266$ 附近。\Cref{fig:bernstein} 显示了 $B_nf_0$ 在尖点附近的
“圆角化”现象：多项式逼近光滑地绕过了 $f_0$ 的不可导点，这正是误差集中于 $x=\frac12$ 的原因。

\begin{table}[htbp]
\centering
\caption{$e_n=B_nf_0(\frac12)$ 的精确值与理论界的比较}\label{tab:sharp}
\begin{tabular}{@{}r S[table-format=1.6] S[table-format=1.6] S[table-format=1.6] S[table-format=1.4]@{}}
\toprule
{$n$} & {$e_n$} & {$\sqrt{n}\,e_n$} & {$\frac32 n^{-1/2}$} & {$e_n\big/\bigl(\frac32n^{-1/2}\bigr)$} \\
\midrule
2    & 0.250000 & 0.353553 & 1.060660 & 0.2357 \\
4    & 0.187500 & 0.375000 & 0.750000 & 0.2500 \\
8    & 0.136719 & 0.386699 & 0.530330 & 0.2578 \\
16   & 0.098190 & 0.392761 & 0.375000 & 0.2618 \\
32   & 0.069975 & 0.395838 & 0.265165 & 0.2639 \\
64   & 0.049673 & 0.397387 & 0.187500 & 0.2649 \\
128  & 0.035193 & 0.398164 & 0.132583 & 0.2654 \\
256  & 0.024910 & 0.398553 & 0.093750 & 0.2657 \\
1024 & 0.012464 & 0.398845 & 0.046875 & 0.2659 \\
\midrule
{$\infty$} & {---} & 0.398942 & {---} & 0.2660 \\
\bottomrule
\end{tabular}
\end{table}

\begin{figure}[htbp]
\centering
\begin{tikzpicture}
\begin{axis}[mojing box, width=.72\textwidth, height=.42\textwidth,
    xmin=0, xmax=1, ymin=0, ymax=.52, xtick={0,.25,.5,.75,1}, ytick={0,.1,.2,.3,.4,.5},
    xlabel={$x$}, grid=major, legend pos=north east, legend columns=2,
    yticklabel style={/pgf/number format/fixed}]
  \addplot[mjInk, line width=.7pt] table[x=x, y=f] {data/bernstein-abs.dat};
  \addlegendentry{$f_0$}
  \addplot[mjGray, dashed] table[x=x, y=B2] {data/bernstein-abs.dat};
  \addlegendentry{$B_2f_0$}
  \addplot[mjSecond] table[x=x, y=B6] {data/bernstein-abs.dat};
  \addlegendentry{$B_6f_0$}
  \addplot[mjAccent] table[x=x, y=B20] {data/bernstein-abs.dat};
  \addlegendentry{$B_{20}f_0$}
\end{axis}
\end{tikzpicture}
\caption{$f_0(x)=\abs{x-\frac12}$ 及其 Bernstein 多项式。$B_nf_0\ge f_0$，数值上最大偏差位于尖点 $x=\frac12$}
\label{fig:bernstein}
\end{figure}

% =============================================================================
\section{Voronovskaya 公式的定量形式}
% =============================================================================

\begin{theorem}\label{thm:voronovskaya}
设 $f\in C^2[0,1]$，则对一切 $n\ge1$ 与 $x\in[0,1]$，
\begin{equation}\label{eq:voronovskaya}
  \Bigl\lvert n\bigl(B_nf(x)-f(x)\bigr)-\frac{x(1-x)}{2}f''(x)\Bigr\rvert
  \le\frac{7}{32}\,\omega\Bigl(f'',\frac{1}{\sqrt n}\Bigr).
\end{equation}
特别地，$n(B_nf-f)$ 在 $[0,1]$ 上一致收敛于 $\frac12x(1-x)f''$。
\end{theorem}

\begin{proof}
由带 Lagrange 余项的 Taylor 公式，$f(x+t)=f(x)+f'(x)t+\frac12f''(x)t^2+r(x,t)\,t^2$，其中
$r(x,t)=\frac12\bigl(f''(\xi)-f''(x)\bigr)$，$\xi$ 介于 $x$ 与 $x+t$ 之间。由 \cref{lem:modulus}，对任意 $\delta>0$，
\[
  \abs{r(x,t)}\le\frac12\,\omega(f'',\abs{t})\le\frac12\Bigl(1+\frac{t^2}{\delta^2}\Bigr)\omega(f'',\delta).
\]
以 $t=T_n$ 代入并取期望，利用 \cref{lem:moments} 中的 $\mathbb{E}\,T_n=0$ 与 $\mathbb{E}\,T_n^2=\frac{x(1-x)}{n}$，得
\[
  \Bigl\lvert n\bigl(B_nf(x)-f(x)\bigr)-\frac{x(1-x)}{2}f''(x)\Bigr\rvert
  =n\,\abs[\big]{\mathbb{E}\,r(x,T_n)T_n^2}
  \le\frac n2\,\omega(f'',\delta)\Bigl(\mathbb{E}\,T_n^2+\frac{\mathbb{E}\,T_n^4}{\delta^2}\Bigr).
\]
取 $\delta=n^{-1/2}$，并用 $\mathbb{E}\,T_n^2\le\frac{1}{4n}$，$\mathbb{E}\,T_n^4<\frac{3}{16n^2}$，右端不超过
$\frac12\bigl(\frac14+\frac{3}{16}\bigr)\omega(f'',n^{-1/2})=\frac{7}{32}\,\omega(f'',n^{-1/2})$。
\end{proof}

\begin{remark}
当 $f''$ 满足 Lipschitz 条件时，\cref{eq:voronovskaya} 给出 $O(n^{-1/2})$ 的误差；对更光滑的 $f$，
误差实际上是 $O(n^{-1})$。以 $f(x)=\sin\pi x$ 为例，此时极限函数为
$-\frac{\pi^2}{2}x(1-x)\sin\pi x$，在 $x=\frac12$ 处等于 $-\frac{\pi^2}{8}$。\Cref{tab:voronovskaya}
中 $n\cdot d_n$ 趋于常数，说明误差 $d_n$ 确为 $O(n^{-1})$；\cref{fig:voronovskaya} 则显示了
$n(B_nf-f)$ 在整个区间上向极限函数的一致逼近。
\end{remark}

\begin{table}[htbp]
\centering
\caption{$f(x)=\sin\pi x$ 时 $n(B_nf-f)(\frac12)$ 的数值，极限为 $-\frac{\pi^2}{8}=-1.233701$}
\label{tab:voronovskaya}
\begin{tabular}{@{}r S[table-format=-1.6] S[table-format=1.6] S[table-format=1.4]@{}}
\toprule
{$n$} & {$n\bigl(B_nf-f\bigr)(\frac12)$} & {$d_n$} & {$n\cdot d_n$} \\
\midrule
5   & -1.109534 & 0.124167 & 0.6208 \\
10  & -1.165148 & 0.068552 & 0.6855 \\
20  & -1.197615 & 0.036086 & 0.7217 \\
40  & -1.215177 & 0.018524 & 0.7409 \\
80  & -1.224315 & 0.009386 & 0.7509 \\
160 & -1.228976 & 0.004724 & 0.7559 \\
320 & -1.231330 & 0.002370 & 0.7585 \\
\bottomrule
\end{tabular}
\par\smallskip
{\footnotesize 注：$d_n=n(B_nf-f)(\frac12)+\frac{\pi^2}{8}$。}
\end{table}

\begin{figure}[htbp]
\centering
\begin{tikzpicture}
\begin{axis}[mojing box, width=.72\textwidth, height=.4\textwidth,
    xmin=0, xmax=1, ymin=-1.35, ymax=.08, xtick={0,.25,.5,.75,1},
    xlabel={$x$}, grid=major, legend columns=4,
    legend style={at={(.5,.97)}, anchor=north},
    yticklabel style={/pgf/number format/fixed}]
  \addplot[mjInk, line width=1pt] table[x=x, y=limit] {data/voronovskaya.dat};
  \addlegendentry{$\frac12x(1-x)f''(x)$}
  \addplot[mjGray, dashed] table[x=x, y=n5] {data/voronovskaya.dat};
  \addlegendentry{$n=5$}
  \addplot[mjSecond] table[x=x, y=n20] {data/voronovskaya.dat};
  \addlegendentry{$n=20$}
  \addplot[mjAccent] table[x=x, y=n80] {data/voronovskaya.dat};
  \addlegendentry{$n=80$}
\end{axis}
\end{tikzpicture}
\caption{$f(x)=\sin\pi x$ 时 $n\bigl(B_nf(x)-f(x)\bigr)$ 与 Voronovskaya 极限的比较}
\label{fig:voronovskaya}
\end{figure}

% =============================================================================
\section{结语}
% =============================================================================

本文的论证只用到了 $B_n$ 是保正线性算子、并且精确再生 $1$、$x$ 两个函数这一事实，以及
二项分布的二阶、四阶矩。因此，同样的方法适用于一切由概率分布生成的保正算子，例如
Sz\'asz--Mirakyan 算子与 Baskakov 算子；这与 Korovkin 定理\cite{korovkin1953}的精神是一致的：
保正线性算子的收敛性由它在少数几个检验函数上的表现决定。

\section*{致谢}
作者感谢审稿人提出的宝贵意见。

\printreferences

\appendix
\section{\cref{eq:exact} 的另一证明}

记 $p_k=\binom{2m}{k}4^{-m}$。由 $\sum_kp_k=1$ 及对称性，$\sum_{k>m}p_k=\frac12(1-p_m)$；由
$kp_k=m\cdot\binom{2m-1}{k-1}2^{-(2m-1)}$，有
$\sum_{k>m}kp_k=m\sum_{j\ge m}\binom{2m-1}{j}2^{-(2m-1)}=\frac m2$。于是
\[
  \sum_{k>m}(k-m)p_k=\frac m2-\frac m2(1-p_m)=\frac m2\,p_m ,
\]
这与正文中的计算一致。

\end{document}
